\documentclass{beamer}
% \usepackage{beamerthemesplit} // Activate for custom appearance
\usepackage{beamerthemesintef} % comment this line + everything with '%%%' if you don't have the theme installed

\titlebackground*{assets/background} %%%
\title{MPC-PiNNs technique applied to Quadrotors}
\subtitle{Elective in Robotics [UAV] Course Project}
\course{Master's Degree in Artificial Intelligence and Robotics} %%%
\author{{Authors}} %%%

\IDnumber{matricole} %%%
\date{Academic Year 2024/2025}

\begin{document}
\maketitle
\section{Introduction}
\section{Models}
\begin{frame}{Unicycle}
    \begin{align}
        \begin{cases}
            \dot{x} = v \cos(\theta) \\
            \dot{y} = v \sin(\theta) \\
            \dot{\theta} = w \\
        \end{cases}
    \end{align}
\end{frame}

\begin{frame}{Quadorotor}
    
\end{frame}
\section{Model Predictive Control}
\section{Physical Informed Neural Networks}
\begin{frame}{PiNNs}
    Physical Informed Neural
Networks (PiNNs) propose an alternative approach where the solution is approxi-
mated by a neural network that is trained not only on data but also by enforcing
the underlying physical laws via the differential equations.

\begin{align}
    \mathcal{L}_{\text{data}}(\theta) &= \frac{1}{N_u} \sum_{i=1}^{N_u} \left| u_{\theta}(\mathbf{x}_i, t_i) - u^{\text{obs}}_i \right|^2 \\
    \mathcal{L}_{\text{phys}}(\theta) &= \frac{1}{N_r} \sum_{j=1}^{N_r} \left| \mathcal{N}[u_{\theta}](\mathbf{x}_j, t_j) \right|^2,
\end{align}

\end{frame}
\section{Results}
\section{Conclusions}
% \backmatter add this later
\end{document}